% ================================================================================
% PEÇA 140: NOTA TÉCNICA DE ATUALIZAÇÃO - POPULAÇÃO DE DOCUMENTOS DO ACHADO I
% ================================================================================
% VERSÃO: 1.0
% DATA: 18 de dezembro de 2025
% ACHADO: I - Inadequações nos Documentos Comprobatórios de Elegibilidade
% REFERÊNCIA: Atualização dos valores constantes na Peça 103
% ================================================================================

\documentclass[11pt,a4paper]{article}

% ========== PACOTES PADRÃO ==========
\usepackage[utf8]{inputenc}
\usepackage[brazil]{babel}
\usepackage[T1]{fontenc}
\usepackage{lmodern}
\usepackage[margin=2.5cm]{geometry}
\usepackage{graphicx}
\usepackage{float}
\usepackage{longtable}
\usepackage{booktabs}
\usepackage{xcolor}
\usepackage{colortbl}
\usepackage{fancyhdr}
\usepackage{amsmath, amssymb}
\usepackage{enumitem}
\usepackage[hidelinks]{hyperref}
\usepackage[breakable]{tcolorbox}
\usepackage{array}

% ========== CORES TCU OFICIAIS ==========
\definecolor{tcured}{RGB}{204,0,0}      % Vermelho TCU #CC0000
\definecolor{tcublue}{RGB}{0,51,102}    % Azul TCU #003366
\definecolor{tcugreen}{RGB}{0,153,76}   % Verde TCU #006633
\definecolor{tcuyellow}{RGB}{255,204,0} % Amarelo TCU #FFCC00
\definecolor{tcuorange}{RGB}{255,153,0} % Laranja TCU #FF9900

% ========== CONFIGURAÇÃO DE CABEÇALHO E RODAPÉ ==========
\pagestyle{fancy}
\fancyhead[L]{Nota Técnica de Atualização - Peça 140}
\fancyhead[R]{TCU/AudTI}
\fancyfoot[C]{\thepage}
\setlength{\headheight}{14pt}

% ========== CONFIGURAÇÃO DE HYPERLINKS ==========
\hypersetup{
    colorlinks=true,
    linkcolor=black,
    filecolor=tcured,
    urlcolor=tcublue,
    citecolor=tcugreen
}

% ========== ESTILOS TCOLORBOX ==========
\newtcolorbox{destaque}[1][]{
    colback=tcuyellow!10,
    colframe=tcuorange,
    fonttitle=\bfseries,
    title=#1,
    breakable
}

\newtcolorbox{importante}[1][]{
    colback=tcured!5,
    colframe=tcured,
    fonttitle=\bfseries,
    title=#1,
    breakable
}

\newtcolorbox{calculo}[1][]{
    colback=tcugreen!5,
    colframe=tcugreen,
    fonttitle=\bfseries,
    title=#1,
    breakable
}

% ========== INFORMAÇÕES DO DOCUMENTO ==========
\title{
    \textbf{NOTA TÉCNICA DE ATUALIZAÇÃO}\\[0.5cm]
    \Large Atualização da População de Documentos e\\
    Projeções Estatísticas do Achado I\\[0.3cm]
    \large (Referência: Peça 103 - Adequação Funcional)
}

\author{
    Equipe de Auditoria Operacional -- TCU/AudTI\\[0.2cm]
    \small TC 011.073/2025-0 -- Auditoria do Sistema CAF
}

\date{18 de dezembro de 2025}

\begin{document}

\maketitle

\begin{abstract}
Esta Nota Técnica documenta a atualização da população total de documentos do Sistema CAF utilizada nas projeções estatísticas do Achado I. Durante a fase de execução da auditoria, identificou-se que a população correta é de \textbf{11.377.318 documentos}, divergindo do valor preliminar de 10.700.156 documentos utilizado na Peça 103. Esta atualização impacta as projeções de documentos problemáticos e os respectivos intervalos de confiança, os quais são recalculados neste documento.
\end{abstract}

\tableofcontents
\newpage

% ================================================================================
\section{Objetivo}
% ================================================================================

Esta Nota Técnica tem por objetivo:

\begin{enumerate}[label=\arabic*)]
    \item Documentar a atualização da população de documentos do Sistema CAF de 10.700.156 para \textbf{11.377.318 documentos};
    \item Explicar a origem da divergência entre os valores;
    \item Apresentar os valores corrigidos para as projeções estatísticas;
    \item Fornecer referência documental para os valores utilizados no relatório de auditoria.
\end{enumerate}

% ================================================================================
\section{Origem da Divergência}
% ================================================================================

\subsection{Valor Preliminar (Peça 103)}

A Peça 103 (Papel de Trabalho de Adequação Funcional) foi elaborada em \textbf{24 de outubro de 2025} utilizando como base a população de documentos disponível naquela data:

\begin{itemize}
    \item \textbf{População utilizada:} 10.700.156 documentos
    \item \textbf{Fonte:} Query inicial sobre a tabela \texttt{S\_DOCUMENTO} do banco CAF
    \item \textbf{Critério:} Contagem de documentos vinculados a cadastros ativos
\end{itemize}

\subsection{Valor Atualizado (População Correta)}

Durante a consolidação das análises e cruzamento com outros papéis de trabalho, identificou-se que a população correta de documentos cadastrados no Sistema CAF é de:

\begin{destaque}[População Correta]
\begin{center}
\textbf{N = 11.377.318 documentos}
\end{center}
\end{destaque}

\subsection{Causa da Divergência}

A diferença de \textbf{677.162 documentos} (6,3\%) entre os valores decorre de:

\begin{enumerate}[label=\alph*)]
    \item \textbf{Critério de contagem:} A query inicial excluía documentos de cadastros em situação ``pendente de análise'', os quais também integram a base documental do sistema;

    \item \textbf{Período de referência:} A extração inicial considerava apenas documentos inseridos até 15/08/2025, enquanto a base completa inclui documentos até a data de extração (02/09/2025);

    \item \textbf{Tabelas auxiliares:} Documentos vinculados a tabelas de relacionamento (ex: \texttt{S\_DOCUMENTO\_UF}) não haviam sido contabilizados na query preliminar.
\end{enumerate}

A população de \textbf{11.377.318 documentos} foi validada por meio de:
\begin{itemize}
    \item Query SQL definitiva executada sobre o dump completo do banco CAF;
    \item Cruzamento com metadados da extração fornecida pelo STI/MAPA (Peça 60);
    \item Verificação independente pela equipe de auditoria.
\end{itemize}

% ================================================================================
\section{Valores Corrigidos}
% ================================================================================

\subsection{Tabela Comparativa: Peça 103 vs. Valores Atualizados}

A tabela a seguir apresenta a comparação entre os valores constantes na Peça 103 e os valores corrigidos que devem ser utilizados como referência no relatório de auditoria:

\begin{table}[H]
\centering
\caption{Comparativo de Valores: Peça 103 vs. Valores Atualizados}
\label{tab:comparativo}
\begin{tabular}{|p{5cm}|>{\centering\arraybackslash}p{3.5cm}|>{\centering\arraybackslash}p{3.5cm}|>{\centering\arraybackslash}p{2.5cm}|}
\hline
\rowcolor{tcublue!20}
\textbf{Parâmetro} & \textbf{Peça 103} & \textbf{Valor Atualizado} & \textbf{Referência} \\
\hline
População total (N) & 10.700.156 & \textbf{11.377.318} & Peça 140, §3 \\
\hline
Proporção problemáticos (p̂) & 27,1\% & 27,1\% & Peça 103, p. 9 \\
\hline
\rowcolor{tcuyellow!20}
\textbf{Projeção pontual} & \textbf{2.900.364} & \textbf{3.083.253} & Peça 140, §4 \\
\hline
IC 99\% - Limite inferior & 2.418.359 & \textbf{2.571.274} & Peça 140, §4 \\
\hline
IC 99\% - Limite superior & 3.382.369 & \textbf{3.595.272} & Peça 140, §4 \\
\hline
\rowcolor{tcuyellow!20}
\textbf{IC 99\% (arredondado)} & \textbf{2,42M a 3,38M} & \textbf{2,57M a 3,60M} & Peça 140, §4 \\
\hline
\end{tabular}
\end{table}

\begin{importante}[Nota]
A \textbf{proporção amostral} (27,1\%) permanece inalterada, pois foi obtida a partir da análise de 646 documentos da amostra. Apenas as \textbf{projeções populacionais} são afetadas pela atualização do denominador (N).
\end{importante}

% ================================================================================
\section{Memória de Cálculo Atualizada}
% ================================================================================

\subsection{Parâmetros Estatísticos}

Os parâmetros estatísticos da análise de adequação funcional são:

\begin{itemize}
    \item \textbf{Tamanho da amostra:} $n = 646$ documentos
    \item \textbf{Documentos problemáticos na amostra:} $x = 175$ (inadequados + parcialmente adequados)
    \item \textbf{Proporção amostral:} $\hat{p} = 175/646 = 0{,}2709$ (27,09\% $\approx$ 27,1\%)
    \item \textbf{Nível de confiança:} 99\% ($z_{0,995} = 2{,}576$)
\end{itemize}

\subsection{Cálculo do Erro Padrão}

\begin{calculo}[Erro Padrão com Correção para População Finita (FPC)]
\begin{align}
SE &= \sqrt{\frac{\hat{p}(1-\hat{p})}{n-1}} \times \sqrt{\frac{N-n}{N-1}} \\[0.3cm]
SE &= \sqrt{\frac{0{,}2709 \times 0{,}7291}{645}} \times \sqrt{\frac{11.377.318 - 646}{11.377.318 - 1}} \\[0.3cm]
SE &= \sqrt{0{,}0003062} \times \sqrt{0{,}9999433} \\[0.3cm]
SE &= 0{,}01750 \times 0{,}99997 \\[0.3cm]
SE &\approx 0{,}0175
\end{align}
\end{calculo}

\textbf{Nota:} O FPC (Finite Population Correction) é praticamente 1,0 dado que $n/N < 0{,}01\%$, mas é mantido por rigor metodológico.

\subsection{Cálculo da Margem de Erro}

\begin{calculo}[Margem de Erro (IC 99\%)]
\begin{align}
ME &= z_{0,995} \times SE \\[0.2cm]
ME &= 2{,}576 \times 0{,}0175 \\[0.2cm]
ME &= 0{,}0451 \approx 4{,}5\text{ p.p.}
\end{align}
\end{calculo}

\subsection{Cálculo do Intervalo de Confiança}

\begin{calculo}[Intervalo de Confiança (99\%)]
\begin{align}
IC_{99\%} &= \hat{p} \pm ME \\[0.2cm]
IC_{99\%} &= 0{,}271 \pm 0{,}045 \\[0.2cm]
IC_{99\%} &= [0{,}226; 0{,}316] \text{ ou } [22{,}6\%; 31{,}6\%]
\end{align}
\end{calculo}

\subsection{Projeção Populacional}

\begin{calculo}[Projeção para a População de 11.377.318 Documentos]
\textbf{Estimativa pontual:}
\begin{align}
\text{Problemáticos} &= N \times \hat{p} \\[0.2cm]
&= 11.377.318 \times 0{,}271 \\[0.2cm]
&= \mathbf{3.083.253} \text{ documentos}
\end{align}

\vspace{0.3cm}
\textbf{Limite inferior (IC 99\%):}
\begin{align}
\text{Mínimo} &= N \times 0{,}226 \\[0.2cm]
&= 11.377.318 \times 0{,}226 \\[0.2cm]
&= \mathbf{2.571.274} \text{ documentos}
\end{align}

\vspace{0.3cm}
\textbf{Limite superior (IC 99\%):}
\begin{align}
\text{Máximo} &= N \times 0{,}316 \\[0.2cm]
&= 11.377.318 \times 0{,}316 \\[0.2cm]
&= \mathbf{3.595.272} \text{ documentos}
\end{align}
\end{calculo}

% ================================================================================
\section{Valores para Uso no Relatório}
% ================================================================================

Com base nos cálculos apresentados, os valores a serem utilizados no relatório de auditoria são:

\begin{destaque}[Valores Oficiais para o Relatório de Auditoria]
\begin{center}
\renewcommand{\arraystretch}{1.5}
\begin{tabular}{|l|r|}
\hline
\rowcolor{tcublue!20}
\textbf{Parâmetro} & \textbf{Valor} \\
\hline
População total de documentos & \textbf{11.377.318} \\
\hline
Proporção de documentos problemáticos & \textbf{27,1\%} \\
\hline
Projeção de documentos problemáticos & \textbf{3,08 milhões} \\
\hline
Intervalo de confiança (99\%) & \textbf{2,57 a 3,60 milhões} \\
\hline
Margem de erro & \textbf{±4,5 p.p.} \\
\hline
\end{tabular}
\end{center}
\end{destaque}

% ================================================================================
\section{Referências Documentais}
% ================================================================================

\begin{table}[H]
\centering
\caption{Rastreabilidade das Referências}
\begin{tabular}{|p{3cm}|p{5cm}|p{6cm}|}
\hline
\rowcolor{tcublue!20}
\textbf{Parâmetro} & \textbf{Peça de Referência} & \textbf{Localização} \\
\hline
População (N) & Peça 140 (esta NT) & Seção 3, Tabela 1 \\
\hline
Proporção (27,1\%) & Peça 103 & Página 9, Tabela 3 \\
\hline
Projeção (3,08M) & Peça 140 (esta NT) & Seção 4, Cálculo de Projeção \\
\hline
IC 99\% (2,57-3,60M) & Peça 140 (esta NT) & Seção 4, Cálculo de IC \\
\hline
Metodologia amostral & Peça 103 & Páginas 5-8 \\
\hline
Exemplos de inadequação & Peça 103 & Páginas 12-14 \\
\hline
\end{tabular}
\end{table}

% ================================================================================
\section{Conclusão}
% ================================================================================

Esta Nota Técnica documenta a atualização da população de documentos do Sistema CAF de 10.700.156 para \textbf{11.377.318 documentos}, bem como as consequentes atualizações nas projeções estatísticas do Achado I.

Os valores atualizados a serem referenciados no relatório de auditoria são:

\begin{itemize}
    \item \textbf{População:} 11.377.318 documentos (ou ``11,4 milhões'')
    \item \textbf{Documentos problemáticos:} 3,08 milhões (IC 99\%: 2,57 a 3,60 milhões)
\end{itemize}

A metodologia de análise, a proporção amostral (27,1\%) e os exemplos de inadequação documentados na Peça 103 permanecem válidos e devem continuar sendo referenciados para esses elementos.

\vspace{1cm}

\begin{flushright}
\textbf{Equipe de Auditoria Operacional}\\
TCU/AudTI\\
18 de dezembro de 2025
\end{flushright}

\end{document}
